\documentclass[a4paper, 12pt]{article}

\usepackage{utils}

\renewcommand*{\today}{24 February 2026}

\begin{document}

\hotbox{Probabilités 1}{CM 6}{\today}

\section{Variables aléatoires réelles}

\begin{definition}
    Soit $(\Omega, \mathcal{T}, P)$ un espace probabilisé. On appelle variable aléatoire
    continue une application $X : \Omega \to \R$ telle que
    $$
    \forall x \in \R, -\infty \leq a < b \leq +\infty, X^{-1}(]a, b]) \in \mathcal{T}
    $$
\end{definition}

\begin{definition}
    Soit $X$ une variable aléatoire continue. Soit, pour $a \in \R$ l'événement
    $$
    \{X \leq a\} = X^{-1}(]-\infty, a])
    $$
    On appelle \textbf{fonction de répartition} de $X$ l'application
    $$
    \funcdef{F_X}{\R}{[0, 1]}{a}{P(X \leq a)}
    $$
\end{definition}

\begin{proposition}{}{}
    Soit $X$ une variable aléatoire continue sur un espace probabilisé $(\Omega, \mathcal{T}, P)$.
    Sa fonction de répartition $F_X$ vérifie les propriétés suivantes :
    \begin{enumerate}
        \item $\forall a < b, P(a < X \leq b) = F_X(b) - F_X(a)$
        \item $F_X$ est croissante
        \item $F_X$ possède une limite à gauche et à droite en tout point de $\R$
        \item $F_X$ est continue à droite en tout point de $\R$
        \item $\lim\limits_{x \to -\infty} F_X(x) = 0$ et $\lim\limits_{x \to +\infty} F_X(x) = 1$
    \end{enumerate}
\end{proposition}

\begin{demonstration}
    \begin{enumerate}
        \item On a
        $$
        X(\omega) \in ]-\infty, b] \iff X(\omega) \leq a \text{ ou } a < X(\omega) \leq b
        $$
        c'est à dire
        $$
        \{ X \leq b \} = \{ X \leq a \} \cup \{ a < X \leq b \}
        $$
        c'est une réunion d'événements incompatibles donc
        $$
        P(X \leq b) = P(X \leq a) + P(a < X \leq b)
        $$
        c'est à dire
        $$
        F_X(b) = F_X(a) + P(a < X \leq b)
        $$
        d'où le résultat.
        \item d'après 1. on a $F_X(b) = F_X(a) + P(a < X \leq b) \geq F_X(a)$ pour tout $a < b$.
        \item $F_X$ est monotone croissante donc elle admet une limite à gauche et à droite en tout point de $\R$.
        \item Soit $x \in \R$ et $(x_n)_{n \geq 1}$ une suite strictement décroissante tel que
        \begin{hotwarn}
            À faire, trompé d'endroit
        \end{hotwarn}
        \begin{enumerate}[label=\roman*.]
            \item $x_{n+1} < x_n, \forall n \geq 1$
            \item $x_n > x, \forall n \geq 1$
            \item $x_n \underset{n \to \infty}{\to} x$
        \end{enumerate}
        On pose $B_n = \{x < X \leq x_n\}$, on a $P(B_n) = F_X(x_n) - F_X(x)$
        on veut montrer que $P(B_n) \underset{n \to \infty}{\to} 0$.
        Comme $x_{n+1} < x_n, \forall n \geq 1$ on a $x < X \leq x_{n+1} \implies x < X \leq x_n$
        d'où $B_{n+1} \subset B_n$.
    \end{enumerate}
\end{demonstration}

\begin{lemme}{}{}
    Soit un espace probabilisé $(\Omega, \mathcal{T}, P)$ et $(B_n)_{n \geq 0}$ une famille d'événements telle que
    \begin{enumerate}
        \item $B_{n+1} \subset B_n, \forall n \geq 0$
        \item $\bigcap\limits_{n \geq 0} B_n = \emptyset$
    \end{enumerate}
    Alors
    $$
    \lim_{n \to \infty} P(B_n) = 0
    $$
\end{lemme}

\begin{theoreme}{}{}
    Soit $X$ une variable aléatoire continue sur un espace probabilisé $(\Omega, \mathcal{T}, P)$.
    L'ensemble des points de discontinuité de sa fonction de répartition $F_X$ est au plus dénombrable.
\end{theoreme}

\begin{remarque}
    Le théorème précédent est vrai pour toute fonction monotone et s'appelle le \textbf{théorème de Froda}.
\end{remarque}

\begin{demonstration}
    Soit $X$ une variable aléatoire continue et une fonction de répartition $F_X$ et $E \subset \R$.
    L'ensemble des points de discontinuité de $F_X$ est
    $$
    E = \{x \in \R \mid F_X(x^-) < F_X(x^+)\}
    $$
    où $F_X(x^-) = \lim\limits_{t \to x^-} F_X(t)$ et $F_X(x^+) = \lim\limits_{t \to x^+} F_X(t)$.
    Soit $x \in E$, il existe $q_x \in ]F_X(x^-), F_X(x^+)[ \cap \Q$.
    On définit une application $\funcdef{\varphi}{E}{\Q}{x}{q_x}$.
    Montrons que $\varphi$ est injective, soit $x, y \in E$ tels que $y > x$, on a $F_X(x^-) < F_X(x^+) \leq F_X(y^-) < F_X(y^+) \implies ]F_X(x^-), F_X(x^+)[\cap]F_X(y^-), F_X(y^+)[ = \emptyset$.
    par définition $\varphi(x) < \varphi(y)$, c'est à dire $\varphi$ est injective. L'ensemble $E$ est au plus dénombrable.
\end{demonstration}

\begin{definition}
    Soit $X$ une variable aléatoire continue sur un espace probabilisé $(\Omega, \mathcal{T}, P)$.
    Supposons que sa fonction de répartition $F_X$ est continue de classe $C^1$ par morceaux.
    On appelle \textbf{densité de probabilité} de $X$ et la dérivée de $F_X$ calculée sur tous les intervalles
    où elle est définie. on la note $f_X$.
\end{definition}

\begin{remarque}
    Au points où $F_X$ n'est pas différentiable, la densité $f_X$ est discontinue mais possède une limite à gauche et à droite.
\end{remarque}

\begin{proposition}{}{}
    Soit $X$ une variable aléatoire continue de densité $f_X$ et de fonction de répartition $F_X$. On a
    \begin{enumerate}
        \item $\forall x \in \R, f_X(x) \geq 0$
        \item $f_X$ est intégrable sur $\R$, d'intégrale égale à 1
        \item pour tout intervalle $]a, b] \subset \R$ on a
        $$
        P(a < X \leq b) = \int_a^b f_X(x)dx
        $$
    \end{enumerate}
\end{proposition}

\begin{demonstration}
    \begin{enumerate}
        \item $F_X$ est croissante, sa dérivée $f_X$ est positive ou nulle.
        \item 
        \begin{align*}
            \int_a^b f_X(x)dx = [F_X(x)]_a^b &= F_X(b) - F_X(a) \\
            &= P(X \in ]a, b])
        \end{align*}
        \item
        $$
        \int_{-M}^M f_X(x)dx = F_X(M) - F_X(-M) \underset{M \to +\infty}{\to} 1
        $$
    \end{enumerate}
\end{demonstration}

\begin{corollaire}{}{}
    Soit $X$ une variable aléatoire continue de densité $f_X$ et de fonction de répartition $F_X$. Alors
    \begin{enumerate}
        \item $\lim\limits_{x \to +\infty} f_X(x) = 0$
        \item Soit $A = ]a, b]$ un intervalle de $\R$. Alors
        $$
        P(X \in A) = 0 \iff f_X|_A \equiv 0
        $$
    \end{enumerate}
\end{corollaire}

\subsection{Espérance et variance d'une variable aléatoire continue}

\begin{definition}
    Soit $X$ une variable aléatoire de densité $f_X$.
    \begin{itemize}
        \item Lorsque $|x| f_X(x)$ est intégrable sur $\R$, on dit
        que $X$ possède \textbf{une espérance}, qu'on définit comme
        $$
        E(X) = \int_{-\infty}^{+\infty} x f_X(x) dx
        $$
        \item De même si $\varphi: \R \to \R$ est une application telle que
        $|\varphi(x)| f_X(x)$ est intégrable sur $\R$, on définit
        $$
        E(\varphi(X)) = \int_{-\infty}^{+\infty} \varphi(x) f_X(x) dx
        $$
    \end{itemize}
\end{definition}

\begin{definition}
    Si $x^2 f_X(x)$ est intégrable sur $\R$, on dit que $X$ possède \textbf{une variance} et on la définit par
    $$
    Var(X) = E((X - \mu)^2) = \int_{-\infty}^{+\infty} (x - \mu)^2 f_X(x) dx
    $$
    avec $\mu = E(X)$.
    On a alors la \textbf{formule de König-Huygens}
    $$
    Var(X) = E(X^2) - (E(X))^2
    $$
\end{definition}

\begin{remarque}
    Si $E(X^2)$ existe, alors $E(X)$ existe et $Var(X)$ existe.
\end{remarque}

\begin{theoreme}{}{}
    Soit $X$ et $Y$ des variables aléatoires continues possèdant une densité, une espérance et une variance. Alors
    \begin{enumerate}
        \item $\forall (a, b) \in \R^2, E(aX + b) = aE(X) + b$
        \item $Var(aX + b) = a^2 Var(X)$
        \item $E(X + Y) = E(X) + E(Y)$
    \end{enumerate}
\end{theoreme}

\end{document}